\chapter{Implémentation}
\label{ch:implementation}

\section{Construction du dataset d'entraînement}

\subsection{Le problème des données disponibles}

Trouver un dataset adapté à ce projet
n'a pas été simple. Le système requiert
la détection fiable de six classes dans
des scènes de surveillance d'intersection
urbaine — bus, car, emergency\_vehicle,
motorcycle, person et truck — avec un
angle de caméra fixe, entre 30 et 60
degrés par rapport au sol, comme dans
les installations réelles aux carrefours.

En explorant les datasets publics disponibles,
trois problèmes se sont posés successivement.

Le premier est l'absence totale de données
marocaines. Aucun dataset public ne représente
le contexte urbain marocain avec nos six
classes. C'est une limite documentée dans
la littérature sur les systèmes de trafic
en Afrique — les données existent surtout
pour l'Europe, l'Asie et l'Amérique du Nord.
Cela signifiait qu'un modèle entraîné sur
des données existantes pourrait ne pas
généraliser correctement aux conditions
visuelles locales : luminosité, densité
de deux-roues, type de véhicules.

Le deuxième problème concerne l'angle de
vue. VisDrone est l'un des datasets les plus
utilisés pour la détection d'objets en milieu
urbain. Il contient de nombreux objets petits,
ce qui semblait pertinent pour notre cas.
Mais après analyse de sa structure et de ses
exemples, il est apparu clairement que ses
images sont prises depuis un drone en vue
quasi-verticale, à 90 degrés. Les objets y
sont minuscules et indifférenciés — les scores
de rappel sur person et motorcycle y sont
respectivement de 36.5\% et 29.6\%, ce qui
est insuffisant pour notre usage. L'angle
d'une caméra de surveillance fixe à
45 degrés est fondamentalement différent :
les véhicules sont vus de face ou de
trois-quarts, avec une perspective qui
change la forme des boîtes englobantes.
VisDrone a donc été écarté à l'analyse,
avant tout entraînement.

Le troisième problème vient de la géométrie
des intersections. La majorité des datasets
disponibles provient d'intersections à double
voie, ce qui rend difficile la délimitation
claire des zones directionnelles Nord, Sud,
Est, Ouest. Sans zones bien définies, le
comptage par zone — qui alimente directement
le module M5 — devient imprécis ou impossible.

Face à ces trois contraintes, la seule
solution était de construire un dataset
hybride à partir de sources sélectionnées
pour leur compatibilité avec nos conditions.

\subsection{Sources retenues}

Quatre sources ont été combinées, chacune
choisie pour combler un manque spécifique :

\begin{table}[H]
\centering
\caption{Sources du dataset final}
\begin{tabularx}{\textwidth}{lXl}
\toprule
\textbf{Source} &
\textbf{Contenu et rôle} &
\textbf{Classes} \\
\midrule
Bellevue Traffic Dataset &
Vidéos d'intersections urbaines
filmées par caméras fixes, angle
de surveillance réel. Source
principale de données véhicules. &
Toutes \\

Lqliaa — Agadir &
Vidéos filmées localement avec
un drone DJI Mini 2 depuis une
hauteur contrôlée, annotées pour
représenter le contexte marocain. &
Toutes \\

Emergency Vehicles &
Images de véhicules d'urgence
en circulation, pour compenser
la rareté de cette classe. &
emergency\_vehicle \\

CityPersons &
Piétons en milieu urbain,
18 villes, tailles variées.
Ajouté pour renforcer la classe
person, doublement critique pour
la décision MDP et l'anonymisation. &
person \\
\bottomrule
\end{tabularx}
\label{tab:dataset_sources}
\end{table}

L'ordre des classes respecte la convention
Roboflow utilisée lors de l'annotation :

\begin{center}
\texttt{bus (0) · car (1) ·
emergency\_vehicle (2) ·
motorcycle (3) · person (4) ·
truck (5)}
\end{center}

\begin{figure}[H]
\centering
\begin{subfigure}{0.48\textwidth}
  \centering
  \includegraphics[width=\textwidth]
    {dataset_bellevue}
  \caption{Bellevue — intersection
           caméra fixe}
\end{subfigure}
\hfill
\begin{subfigure}{0.48\textwidth}
  \centering
  \includegraphics[width=\textwidth]
    {dataset_lqliaa}
  \caption{Lqliaa — contexte marocain}
\end{subfigure}
\caption{Exemples de scènes du dataset final
         \textit{(Source : auteur)}}
\label{fig:dataset_sources_visual}
\end{figure}

Bellevue apporte la densité de véhicules
et la géométrie d'intersection compatible
avec nos zones. Les images de Lqliaa
ajoutent le contexte visuel marocain :
luminosité différente, types de véhicules
locaux, densité de motos. Les deux sources
sont complémentaires.

\subsection{Méthodologie de construction
            en trois étapes}

La construction a suivi une progression
en trois étapes, chacune répondant à
un problème identifié sur les données
précédentes.

\subsubsection{Étape 1 — Split sans fuite
               de données}

Les images étant extraites de vidéos,
des frames consécutives d'une même
séquence peuvent se retrouver à la fois
dans le train et dans le test. Si c'est
le cas, le modèle a vu pendant
l'entraînement des frames quasi-identiques
à celles qu'on utilise pour l'évaluer —
les métriques semblent bonnes, mais elles
ne reflètent pas la vraie capacité de
généralisation.

Pour éviter ça, le split a été fait par
groupe de source : toutes les frames
extraites d'une même vidéo vont dans
le même split. La vérification confirme
\texttt{0 group(s) found in multiple splits}.

\begin{table}[H]
\centering
\caption{Distribution du dataset — Étape 1}
\small
\begin{tabularx}{\textwidth}{lrrrrrrr}
\toprule
\textbf{Split} &
\textbf{Images} &
\textbf{bus} &
\textbf{car} &
\textbf{emerg.} &
\textbf{moto} &
\textbf{person} &
\textbf{truck} \\
\midrule
Train & 2716 & 407 & 19268 & 286 &
4358 & 2075 & 1341 \\
Val   & 779  & 117 & 5486  & 83  &
1241 & 507  & 384  \\
Test  & 393  & 62  & 2742  & 42  &
622  & 309  & 200  \\
\midrule
\textbf{Total} & \textbf{3888} &
586 & 27496 & 411 &
6221 & 2891 & 1925 \\
\bottomrule
\end{tabularx}
\label{tab:step1}
\end{table}

\subsubsection{Étape 2 — Renforcement
               de la classe person}

Avec 2 075 instances person en train,
la classe était sous-représentée par
rapport à car (19 268 instances). C'est
un problème doublement critique dans
notre système : la détection des piétons
alimente à la fois la décision MDP et
l'anonymisation Privacy-by-Design. Un
piéton non détecté n'est pas anonymisé
— ce qui constitue une violation de la
Loi 09-08.

CityPersons a été intégré exclusivement
dans l'ensemble d'entraînement. Ajouter
des données de renforcement dans val
ou test biaiserait l'évaluation en
faveur des classes renforcées.

\begin{table}[H]
\centering
\caption{Ajout CityPersons — Étape 2}
\begin{tabularx}{\textwidth}{lcc}
\toprule
\textbf{Métrique} &
\textbf{Avant} &
\textbf{Après} \\
\midrule
Images train & 2716 & 5058 \\
Instances person (train) &
2075 & 20652 \\
Villes représentées & — & 18 \\
Val/Test & inchangés & inchangés \\
\bottomrule
\end{tabularx}
\label{tab:step2}
\end{table}

\subsubsection{Étape 3 — Crops pour
               tiny persons}

Dans les vidéos de surveillance
d'intersection, les piétons apparaissent
souvent très petits — moins de 0.1\%
de la surface totale de l'image quand
ils sont éloignés de la caméra. YOLO
a tendance à ignorer ces petites
détections, ce qui pose exactement
le même problème que l'étape précédente
pour l'anonymisation.

Pour traiter ce cas, un script Python
(\texttt{create\_tiny\_person\_step3.py})
a généré automatiquement 531 crops
centrés sur les boîtes person de surface
normalisée $\leq 0.10\%$. Pour chaque
crop, toutes les classes visibles dans
la région recadrée ont été conservées
avec recalcul des coordonnées YOLO.
Ces crops ont été ajoutés uniquement
au train.

\begin{table}[H]
\centering
\caption{Distribution finale — Étape 3}
\small
\begin{tabularx}{\textwidth}{lrrrrrrr}
\toprule
\textbf{Split} &
\textbf{Images} &
\textbf{bus} &
\textbf{car} &
\textbf{emerg.} &
\textbf{moto} &
\textbf{person} &
\textbf{truck} \\
\midrule
Train & 5589 & 468 & 23071 & 286 &
4387 & 22280 & 1531 \\
Val   & 779  & 117 & 5486  & 83  &
1241 & 507   & 384  \\
Test  & 393  & 62  & 2742  & 42  &
622  & 309   & 200  \\
\midrule
\textbf{Total} & \textbf{6761} &
647 & 31299 & 411 &
6250 & 23096 & 2115 \\
\bottomrule
\end{tabularx}
\label{tab:step3_final}
\end{table}

La distribution par taille de boîte
sur la classe person montre que le
dataset final couvre toutes les
tailles rencontrées en surveillance
réelle :

\begin{table}[H]
\centering
\caption{Tailles de boîtes person —
         Dataset final}
\small
\begin{tabularx}{\textwidth}{lXXXXX}
\toprule
\textbf{Split} &
\textbf{Tiny} &
\textbf{V. Small} &
\textbf{Small} &
\textbf{Medium} &
\textbf{Large} \\
& \scriptsize{<0.05\%} &
\scriptsize{0.05-0.1\%} &
\scriptsize{0.1-0.25\%} &
\scriptsize{0.25-1\%} &
\scriptsize{$\geq$1\%} \\
\midrule
Train & 6902 & 3707 & 5217 & 4689 & 1765 \\
Val   & 404  & 70   & 31   & 2    & 0    \\
Test  & 248  & 49   & 12   & 0    & 0    \\
\bottomrule
\end{tabularx}
\label{tab:person_sizes}
\end{table}

La majorité des instances person dans
val et test sont de taille tiny ou
very small — ce qui correspond à la
réalité des vidéos de surveillance
d'intersection et explique les limites
de rappel observées sur cette classe.

\subsection{Audit qualité}

Un audit visuel des annotations a été
conduit sur les ensembles val et test
via des contact sheets générés
automatiquement. L'objectif était de
vérifier la cohérence des boîtes
englobantes, notamment sur les petites
instances person, et d'identifier
d'éventuelles erreurs d'annotation
avant l'entraînement final.

\begin{figure}[H]
\centering
\begin{subfigure}{0.49\textwidth}
  \centering
  \includegraphics[width=\textwidth]
    {annotation}
  \caption{Contrôle des annotations}
\end{subfigure}
\hfill
\begin{subfigure}{0.49\textwidth}
  \centering
  \includegraphics[width=\textwidth]
    {class_distribution}
  \caption{Distribution par classe}
\end{subfigure}
\caption{Audit visuel du dataset final
         \textit{(Source : auteur)}}
\label{fig:dataset_audit}
\end{figure}

Cet audit a mis en évidence le déséquilibre
attendu entre car — classe dominante —
et les classes rares comme
emergency\_vehicle, bus et truck. Ce
déséquilibre a été compensé dans
l'entraînement par class weighting.

\section{Benchmark et sélection du modèle}

\subsection{Protocole expérimental}

Le benchmark a été conduit en deux phases
complémentaires. La première phase
d'entraînement s'est faite sur GPU
(Google Colab Pro, NVIDIA A100) pour
mesurer la précision. La seconde phase
de déploiement s'est faite directement
sur Raspberry Pi 5 pour mesurer le débit
réel en conditions embarquées. Ces deux
phases sont nécessaires parce que les
métriques GPU ne préjugent pas du débit
sur CPU ARM — un modèle précis sur GPU
peut être inutilisable sur Pi 5.

Six architectures YOLO ont été évaluées
sur le même dataset Step 3 (6 761 images),
avec les mêmes hyperparamètres de base :

\begin{table}[H]
\centering
\caption{Modèles évalués — Protocole}
\begin{tabularx}{\textwidth}{llccc}
\toprule
\textbf{Modèle} &
\textbf{Variante} &
\textbf{Batch} &
\textbf{Epochs} &
\textbf{ImgSz} \\
\midrule
YOLOv8n  & Nano  & 64 & 100 & 640 \\
YOLOv8s  & Small & 32 & 100 & 640 \\
YOLOv11n & Nano  & 64 & 100 & 640 \\
YOLOv11s & Small & 32 & 100 & 640 \\
YOLO26n  & Nano  & 64 & 100 & 640 \\
YOLO26s  & Small & 32 & 100 & 640 \\
\bottomrule
\end{tabularx}
\label{tab:benchmark_protocol}
\end{table}

Chaque modèle a été entraîné avec
class weighting pour compenser le
déséquilibre des classes, early stopping
et augmentation avancée. Les meilleurs
poids \texttt{best.pt} ont ensuite été
exportés en ONNX pour le déploiement.

\subsection{Résultats GPU — précision}

Le fine-tuning final a porté sur YOLO26n
en deux résolutions successives. On est
parti des poids Step 3 800 déjà stabilisés,
puis on a augmenté la résolution à 960
pixels pour améliorer la détection des
petits objets.

\begin{table}[H]
\centering
\caption{YOLO26n Step 3 — comparaison
         800 vs 960 pixels}
\small
\begin{tabularx}{\textwidth}{Xcccc}
\toprule
\textbf{Modèle} &
\textbf{mAP50} &
\textbf{mAP50-95} &
\textbf{R person} &
\textbf{mAP50 person} \\
\midrule
YOLO26n Step 3 800 &
0.775 & 0.563 & 0.291 & 0.450 \\
\rowcolor{PrimaryBlue!10}
YOLO26n Step 3 960 &
\textbf{0.792} & \textbf{0.579} &
\textbf{0.353} & \textbf{0.512} \\
\bottomrule
\end{tabularx}
\label{tab:gpu_benchmark}
\end{table}

Le gain est visible sur toutes les métriques,
mais surtout sur la classe person : le rappel
passe de 0.291 à 0.353 et la mAP50 person
de 0.450 à 0.512. C'est attendu — une
résolution plus élevée aide directement
pour les petits objets.

Les paramètres du dernier run :

\begin{table}[H]
\centering
\caption{Paramètres du fine-tuning
         YOLO26n Step 3 960}
\small
\begin{tabularx}{\textwidth}{lX}
\toprule
\textbf{Paramètre} & \textbf{Valeur} \\
\midrule
Poids de départ &
\texttt{YOLO26n\_step3\_800\_best.pt} \\
GPU &
NVIDIA A100 40GB — Google Colab Pro \\
Résolution &
960$\times$960 pixels \\
Batch effectif &
26 images (AutoBatch) \\
Optimiseur &
AdamW, lr0=1e-4, cosine decay \\
Augmentation &
mosaic=0.1, HSV, flips horizontaux \\
Durée effective &
$\approx$36 min (2150s),
arrêt epoch 25, meilleur epoch 15 \\
\bottomrule
\end{tabularx}
\label{tab:training_params}
\end{table}

L'entraînement s'est terminé proprement
par early stopping à l'epoch 25, avec
le meilleur modèle sauvegardé à l'epoch 15.
La session Colab complète — montage Drive,
décompression dataset, entraînement, export
et téléchargement des poids — a pris environ
45 minutes à 1 heure.

\begin{figure}[H]
\centering
\includegraphics[width=0.98\textwidth]
  {results_cropped}
\caption{Courbes d'entraînement YOLO26n
         Step 3 960 \textit{(Source : auteur)}}
\label{fig:training_curves}
\end{figure}

\begin{figure}[H]
\centering
\begin{subfigure}{0.49\textwidth}
  \centering
  \includegraphics[width=\textwidth]
    {confusion_matrix_cropped}
  \caption{Matrice de confusion}
\end{subfigure}
\hfill
\begin{subfigure}{0.49\textwidth}
  \centering
  \includegraphics[width=\textwidth]
    {BoxPR_curve_cropped}
  \caption{Courbe précision-rappel}
\end{subfigure}
\caption{Évaluation du modèle final
         \textit{(Source : auteur)}}
\label{fig:eval_plots}
\end{figure}

La matrice de confusion confirme une limite
attendue : la confusion entre car, bus et
truck est la plus fréquente. Ce n'était pas
surprenant — dans les vues CCTV, ces véhicules
sont vus de face ou de trois-quarts avec des
formes proches quand ils sont éloignés,
partiellement occultés ou dans des conditions
de faible luminosité. C'est une limite
inhérente au modèle nano sur des classes
visuellement proches, indépendante du
dataset choisi.

\subsection{Résultats Pi 5 — débit réel}

Le benchmark sur Raspberry Pi 5 a été
conduit avec ONNX Runtime en mode CPU,
sur 200 frames pour les six architectures
de base et 300 frames pour les variantes
Step 3 :

\begin{table}[H]
\centering
\caption{Configuration benchmark Pi 5}
\begin{tabularx}{\textwidth}{lX}
\toprule
\textbf{Paramètre} & \textbf{Valeur} \\
\midrule
Hardware & Raspberry Pi 5, 8GB RAM \\
Runtime  & ONNX Runtime 1.26.0 \\
Provider & CPUExecutionProvider \\
Résolution & 640 (benchmark famille) ;
960 et 800 (Step 3) \\
Frames testées & 200 (famille) ;
300 (Step 3) \\
Vidéo de test & Bellevue NE8th — jour \\
\bottomrule
\end{tabularx}
\label{tab:rpi5_config}
\end{table}

\begin{table}[H]
\centering
\caption{Benchmark complet sur Raspberry Pi 5}
\small
\begin{tabularx}{\textwidth}{llccccc}
\toprule
\textbf{Modèle} &
\textbf{Format} &
\textbf{FPS} &
\textbf{mAP@50} &
\textbf{Emergency} &
\textbf{Person} &
\textbf{Taille} \\
\midrule
\multicolumn{7}{l}{\textit{Variantes Nano —
compatibles temps réel ($\geq$ 5 FPS)}} \\
\midrule
\rowcolor{PrimaryBlue!10}
\textbf{YOLO26n} & \textbf{FP32} &
\textbf{7.6} & \textbf{65.6\%} &
\textbf{70.4\%} & \textbf{41.5\%} &
\textbf{9.8MB} \\
YOLO26n  & INT8 & 6.2 & 64.0\% &
69.0\% & 40.0\% & 2.9MB \\
YOLOv11n & FP32 & 5.9 & 69.2\% &
73.4\% & 44.8\% & 10.6MB \\
YOLOv8n  & INT8 & 5.7 & 68.0\% &
78.0\% & 40.0\% & 3.4MB \\
YOLOv11n & INT8 & 5.7 & 68.0\% &
72.0\% & 43.0\% & 3.0MB \\
YOLOv8n  & FP32 & 5.3 & 69.8\% &
80.1\% & 41.5\% & 12.3MB \\
\midrule
\multicolumn{7}{l}{\textit{Variantes Small —
incompatibles temps réel ($<$ 3 FPS)}} \\
\midrule
YOLO26s  & INT8 & 2.5 & 73.0\% &
85.0\% & 56.0\% & 10.0MB \\
YOLOv11s & INT8 & 2.5 & 72.0\% &
78.0\% & 51.0\% & 9.9MB \\
YOLOv8s  & INT8 & 2.4 & 69.0\% &
77.0\% & 50.0\% & 11.5MB \\
YOLO26s  & FP32 & 2.1 & 74.9\% &
87.1\% & 58.0\% & 38.2MB \\
YOLOv11s & FP32 & 1.9 & 73.5\% &
80.2\% & 52.6\% & 37.9MB \\
YOLOv8s  & FP32 & 1.5 & 71.0\% &
79.1\% & 51.2\% & 44.8MB \\
\bottomrule
\end{tabularx}
\label{tab:full_benchmark}
\end{table}

\subsection{Décision finale}

Deux résultats ressortent clairement
de ce benchmark.

Le premier concerne les variantes small.
Leur incompatibilité avec le temps réel
sur Pi 5 n'était pas totalement
surprenante — elles sont plus coûteuses
en calcul. Mais le benchmark a permis
de la quantifier précisément : aucune
variante small ne dépasse 3 FPS, quelle
que soit la précision numérique du format
d'export. Ce constat ferme définitivement
la porte aux architectures small pour
un déploiement sur CPU ARM.

Le deuxième concerne la quantification
INT8. Elle réduit la taille du modèle
significativement, mais n'apporte pas
de gain FPS systématique et introduit
une confusion accrue entre car, bus et
truck. Le redéploiement final du modèle
Step 3 800 en INT8 dynamique en est
l'illustration : la taille passe de
9.4 MB à 2.8 MB, mais le débit descend
à 2.06 FPS contre 3.94 FPS en FP32.
Le format FP32 est donc retenu.

YOLO26n FP32 est sélectionné pour trois
raisons concrètes. Son débit de 7.6 FPS
sur Pi 5 est le plus élevé parmi les
variantes nano — et la marge au-dessus
du seuil minimal de 5 FPS laisse des
ressources CPU disponibles pour le reste
du pipeline. Son architecture NMS-Free
supprime l'étape de post-traitement qui
consomme du CPU sur les autres modèles.
Et son empreinte de 9.8 MB en ONNX FP32
est la plus légère des trois architectures
nano retenues.

La variante Step 3 960 est retenue comme
modèle de précision. Son redéploiement
complet avec ROI atteint 2.16 FPS sur Pi 5
— ce qui en fait un mode d'analyse plutôt
que de traitement temps réel continu.
Un export 800 pixels depuis les mêmes poids
améliore le débit à 3.94 FPS sans ROI.
Avec \texttt{--vid-stride 2}, ce même
export couvre 7.84 frames source par
seconde en n'inférant qu'une frame sur deux
— le meilleur compromis embarqué obtenu
pour le modèle Step 3.

\begin{table}[H]
\centering
\caption{Récapitulatif des variantes
         Step 3 sur Raspberry Pi 5}
\small
\begin{tabularx}{\textwidth}{Xcccl}
\toprule
\textbf{Variante} &
\textbf{FPS analyse} &
\textbf{FPS source} &
\textbf{Taille} &
\textbf{Statut} \\
\midrule
Step 3 960 + ROI &
2.16 & 2.16 & 9.6MB &
Mode haute précision \\
Step 3 960 sans ROI &
2.51 & 2.51 & 9.6MB &
Limite inférence 960 \\
Step 3 800 sans ROI &
3.94 & 3.94 & 9.4MB &
Export expérimental \\
Step 3 800 stride 2 &
4.28 & 7.84 & 9.4MB &
Meilleur compromis \\
Step 3 800 INT8 &
2.06 & 2.06 & 2.8MB &
Écarté — trop lent \\
\bottomrule
\end{tabularx}
\label{tab:step3_rpi5}
\end{table}

\section{Traçabilité et reproductibilité}

\subsection{Approche adoptée}

MLflow n'a pas piloté tous les
entraînements depuis le début. Il a été
intégré comme couche de documentation
finale, une fois le modèle retenu, pour
centraliser les paramètres, les métriques
et les artefacts de déploiement dans un
format reproductible. L'objectif n'était
pas de construire une infrastructure MLOps
industrielle, mais de rendre les choix
du projet vérifiables et répétables.

Trois niveaux de traçabilité ont été mis
en place.

Le premier niveau concerne les données
et paramètres. Le dataset final est décrit
par \texttt{data.yaml}. Les paramètres
d'entraînement et d'inférence sont
centralisés dans \texttt{params.yaml}.
Les métriques finales sont conservées
dans \path{metrics/final_yolo26n_step3_960.json}
et les résultats terrain dans
\path{results/pipeline_results.json}.

Le deuxième niveau est le suivi MLflow.
Le script \texttt{scripts/log\_mlflow\_final.py}
crée une exécution MLflow locale (base
SQLite \texttt{mlflow.db}) contenant les
paramètres, les métriques et les artefacts
finaux — poids PyTorch, export ONNX,
courbes CSV. L'exécution finale porte
le nom \texttt{YOLO26n-Step3-960-final-6582f915}.

Le troisième niveau est le déploiement.
La commande \texttt{python3 -m pipeline.main}
standardise le lancement du pipeline.
L'export ONNX assure la portabilité vers
le Raspberry Pi 5. Le dashboard Flask
permet de valider le fonctionnement
en conditions réelles.

\begin{figure}[H]
\centering
\includegraphics[width=0.98\textwidth]
  {mlflow_ui_report}
\caption{Interface MLflow — expérience finale
         \texttt{smart-traffic-agadir-yolo26n}
         \textit{(Source : auteur)}}
\label{fig:mlflow_ui}
\end{figure}

\begin{figure}[H]
\centering
\includegraphics[width=0.95\textwidth]
  {mlflow_artifacts_report}
\caption{Artefacts du run final :
         poids PyTorch, ONNX, métriques
         \textit{(Source : auteur)}}
\label{fig:mlflow_artifacts}
\end{figure}

\begin{figure}[H]
\centering
\resizebox{0.94\textwidth}{!}{%
\begin{tikzpicture}[
  font=\small,
  box/.style={
    rounded corners=5pt, very thick,
    minimum width=4.0cm,
    minimum height=1.08cm,
    align=center
  },
  data/.style={box,
    draw=PrimaryBlue, fill=PrimaryBlue!8},
  exp/.style={box,
    draw=AccentGreen, fill=AccentGreen!8},
  dep/.style={box,
    draw=AccentOrange,
    fill=AccentOrange!10},
  lane/.style={
    rounded corners=8pt,
    draw=DarkGray!45,
    fill=LightGray!45,
    inner xsep=0.35cm,
    inner ysep=0.30cm
  },
  arrow/.style={
    -{Latex[length=3mm]},
    line width=1.15pt,
    color=DarkGray},
  note/.style={
    font=\scriptsize, color=DarkGray,
    fill=white, inner sep=2pt}
]

\node[data] (raw) at (0,0)
  {Données brutes\\vidéos + images};
\node[data] (prep) at (5.0,0)
  {Préparation\\splits + crops};
\node[data] (yaml) at (10.0,0)
  {\texttt{data.yaml}\\classes + chemins};

\node[exp] (train) at (10.0,-2.20)
  {Colab A100\\fine-tuning};
\node[exp] (mlflow) at (5.0,-2.20)
  {MLflow local\\paramètres + artefacts};
\node[exp] (metrics) at (0,-2.20)
  {\texttt{metrics/}\\résultats finaux};

\node[dep] (onnx) at (0,-4.40)
  {Export ONNX\\FP32};
\node[dep] (pi) at (5.0,-4.40)
  {Raspberry Pi 5\\benchmark CPU};
\node[dep] (dash) at (10.0,-4.40)
  {Dashboard\\validation};

\begin{scope}[on background layer]
\node[lane,
  fit=(raw)(prep)(yaml),
  label={[font=\bfseries\color{PrimaryBlue}]
    left:Dataset}] {};
\node[lane,
  fit=(train)(mlflow)(metrics),
  label={[font=\bfseries\color{AccentGreen}]
    left:MLOps}] {};
\node[lane,
  fit=(onnx)(pi)(dash),
  label={[font=\bfseries\color{AccentOrange}]
    left:Déploiement}] {};
\end{scope}

\draw[arrow] (raw) --
  node[note, above] {sélection} (prep);
\draw[arrow] (prep) --
  node[note, above] {configuration} (yaml);
\draw[arrow] (yaml) --
  node[note, right] {données} (train);
\draw[arrow] (train) --
  node[note, above] {run} (mlflow);
\draw[arrow] (mlflow) --
  node[note, above] {logs} (metrics);
\draw[arrow] (metrics) --
  node[note, left] {modèle final} (onnx);
\draw[arrow] (onnx) --
  node[note, above] {déploiement} (pi);
\draw[arrow] (pi) --
  node[note, above] {preuves} (dash);

\node[
  draw=DarkGray!50, fill=white,
  rounded corners=4pt, align=center,
  font=\scriptsize, text width=11.3cm,
  below=0.65cm of pi
]{MLflow centralise les paramètres
et artefacts du run final. Le benchmark
Pi 5 fournit les mesures de débit réel.};
\end{tikzpicture}}
\caption{Chaîne de traçabilité du projet
         \textit{(Source : auteur)}}
\label{fig:mlops_pipeline}
\end{figure}

\subsection{Commande de lancement}

Le pipeline final se lance depuis la
racine du projet avec :

\begin{lstlisting}[language=bash,
caption={Lancement du pipeline final}]
python3 -m pipeline.main \
  --video data/videos/ne8th/Bellevue_NE8th.mp4 \
  --model models/downloads/YOLO26n_step3_960_best.onnx \
  --imgsz 960 \
  --conf 0.35 \
  --dashboard \
  --person-roi \
  --vehicle-roi
\end{lstlisting}

Le dashboard est ensuite accessible à
\texttt{http://<ip-raspberry>:5000}.
Les chemins complets des artefacts et
fichiers de configuration sont détaillés
dans le dépôt du projet.